\subsection{Green\textquotesingle s Functions}\label{sec:greens-functions}

In the previous two sections, we have considered fundamental solutions for PDEs with constant coefficients. Such fundamental solutions allow the solution of problems of the form $L(D)u=f$, posed on all of space. In practical applications, however, one does not usually want to solve problems posed on all of space; rather one wants to solve PDEs on some domain, subject to certain conditions on the boundary. Green\textquotesingle s functions are the analogue of fundamental solutions for this situation. They can be found explicitly only in very special cases. Nevertheless, the concept of Green\textquotesingle s functions is useful for theoretical investigations. At present, we do not have the methods available to discuss the existence, uniqueness and regularity of Green\textquotesingle s functions for PDEs, and the discussions in this section will to a large extent be formal.

\subsubsection{Boundary-Value Problems and their Adjoints}\label{subsec:boundary-value-problems-adjoints}

\begin{Definition}
\label{def:formal-adjoint-differential-operator}
Let
\begin{equation}\label{eq:variable-coefficient-differential-operator}
L(\bm{x},D)=\sum_{|\alpha|\leq k}a_\alpha(\bm{x})D^\alpha
\end{equation}
be a differential operator defined on $\Omega$. Then the \textbf{formal adjoint} of $L$ is the operator given by
\begin{equation}\label{eq:formal-adjoint-differential-operator}
L^*(\bm{x},D)u=\sum_{|\alpha|\leq k}(-1)^{|\alpha|}D^\alpha(a_\alpha(\bm{x})u(\bm{x})).
\end{equation}

\end{Definition}

The importance of this definition lies in the fact that
\begin{equation}\label{eq:formal-adjoint-test-function-identity}
(\phi,L(\bm{x},D)\psi)=(L^*(\bm{x},D)\phi,\psi)
\end{equation}
for every $\phi,\psi\in\mathcal{D}(\Omega)$. If the assumption of compact support is removed, then \eqref{eq:formal-adjoint-test-function-identity} no longer holds; instead the integration by parts yields additional terms involving integrals over the boundary $\partial\Omega$. However, these boundary terms vanish if $\phi$ and $\psi$ satisfy certain restrictions on the boundary. We are interested in the case where the order of $L$ is even, $k=2p$ and $p$ linear homogeneous boundary conditions on $\psi$ are given. It is then natural to seek $p$ boundary conditions to be imposed on $\phi$ which would make \eqref{eq:formal-adjoint-test-function-identity} hold. This leads to the notion of an adjoint boundary-value problem.

To make this idea concrete, let us first consider the case of ordinary differential equations. Let
\begin{equation}\label{eq:even-order-ordinary-differential-operator}
L(x,D)u=\sum_{i=0}^{2p}a_i(x)\dfrac{\mathrm{d}^iu}{\mathrm{d}x^i}(x),\quad x\in(a,b).
\end{equation}
We assume that $a_i\in C^i([a,b])$; this guarantees that the coefficients of $L^*$ are continuous. Moreover, we assume that $a_{2p}(x)\ne0$ for $x\in[a,b]$. For any functions $\phi,\psi\in C^{2p}[a,b]$, we compute
\begin{equation}\label{eq:ode-formal-adjoint-boundary-expansion}
\begin{aligned}
&(\phi,L(x,D)\psi)-(L^*(x,D)\phi,\psi)\\
&\quad=\sum_{i=1}^{2p}\sum_{k=0}^{i-1}(-1)^kD^{i-k-1}\psi(b)D^k(a_i\phi)(b)\\
&\qquad-\sum_{i=1}^{2p}\sum_{k=0}^{i-1}(-1)^kD^{i-k-1}\psi(a)D^k(a_i\phi)(a).
\end{aligned}
\end{equation}


The boundary terms can be recast in the form
\begin{equation}\label{eq:ode-adjoint-boundary-matrix-form}
\begin{aligned}
&\sum_{k,l=1}^{2p}A_{kl}(b)D^{k-1}\phi(b)D^{l-1}\psi(b)\\
&\qquad-\sum_{k,l=1}^{2p}A_{kl}(a)D^{k-1}\phi(a)D^{l-1}\psi(a).
\end{aligned}
\end{equation}
Here $A_{kl}$ vanishes for $k+l>2p+1$, and $A_{k(2p+1-k)}=(-1)^{k-1}a_{2p}$. Since $a_{2p}$ was assumed nonzero, this implies that the matrix $\bm{A}$ is nonsingular. Now assume that at the point $b$ we have $p$ linearly independent boundary conditions
\begin{equation}\label{eq:ode-independent-endpoint-boundary-conditions}
\sum_{j=1}^{2p}b_{ij}D^{j-1}\psi(b)=0,\quad i=1,\ldots,p.
\end{equation}
Let $\bm{u}$ denote the $2p$ vector with components $u_i=D^{i-1}\psi(b)$ and let $\bm{v}$ denote the $2p$ vector with components $v_i=D^{i-1}\phi(b)$. Then \eqref{eq:ode-independent-endpoint-boundary-conditions} constrains $\bm{u}$ to a $p$-dimensional subspace $X$ of $\mathbb{R}^{2p}$. The image of $X$ under $\bm{A}(b)$ is a $p$-dimensional subspace $Y$ of $\mathbb{R}^{2p}$. In order to make the first term in \eqref{eq:ode-adjoint-boundary-matrix-form} vanish for every $\psi$ that satisfies \eqref{eq:ode-independent-endpoint-boundary-conditions}, it is necessary and sufficient to have $\bm{v}$ in the orthogonal complement of $Y$. This yields a set of $p$ boundary conditions for $\phi$, which we call the adjoint boundary conditions. An analogous consideration applies at the point $a$.

As an example, let $L\psi=\psi''''$ with the boundary conditions $\psi+\psi^{\prime}=2\psi+\psi^{\prime\prime}=0$ at each endpoint. In this case, $L^*=L$, and \eqref{eq:ode-formal-adjoint-boundary-expansion} specializes to
\begin{equation}\label{eq:fourth-order-ode-green-identity}
\begin{aligned}
\int_a^b\phi(x)\psi''''(x)\,\mathrm{d}x
&=\int_a^b\phi''''(x)\psi(x)\,\mathrm{d}x\\
&\quad+\phi(b)\psi'''(b)-\phi'(b)\psi''(b)\\
&\quad+\phi''(b)\psi'(b)-\phi'''(b)\psi(b)\\
&\quad-\phi(a)\psi'''(a)+\phi'(a)\psi''(a)\\
&\quad-\phi''(a)\psi'(a)+\phi'''(a)\psi(a).
\end{aligned}
\end{equation}
The matrix $\bm{A}$ in \eqref{eq:ode-adjoint-boundary-matrix-form} is
\begin{equation}\label{eq:fourth-order-ode-boundary-matrix}
\bm{A}=\begin{pmatrix}0&0&0&1\\0&0&-1&0\\0&1&0&0\\-1&0&0&0\end{pmatrix},
\end{equation}
and the vector $\bm{u}$ is subject to the conditions $u_1+u_2=2u_1+u_3=0$. A basis for the space $X$ of vectors satisfying these conditions is given by the vectors $(1,-1,-2,0)$ and $(0,0,0,1)$. The images of these vectors under $\bm{A}$ are $(0,2,-1,-1)$ and $(1,0,0,0)$. Thus the vector $\bm{v}$ has to satisfy the conditions $2v_2-v_3-v_4=0$, $v_1=0$, i.e., the adjoint boundary conditions are $\phi=2\phi^{\prime}-\phi^{\prime\prime}-\phi^{\prime\prime\prime}=0$.
